\section{Discussion}
\label{sec:discussion}

\subsection{Key Findings}

\textbf{RLEF's difficulty-scaling advantage is real and statistically supported, but the mechanism differs from prior assumptions.} JT z=+56.10 confirms positive-ordered trend; LCB-Hard shows the largest $\Delta$ (+0.18). This is consistent with \cite{Gehring2024RLEF} and extends it to a reproducible open-source setting.

The generalization void observation has broader implications: as code LLMs scale, the key barrier to hard-benchmark performance may not be dataset coverage but distribution shift between training solutions and evaluation problem formats. RLEF's self-generated feedback mechanism is well-suited to this regime; SFT is not. We emphasize this is a hypothesis suggested by our data, not a mechanism we have directly measured.

\textbf{The reward formulation null result is practically important.} Fraction-of-tests and binary rewards are equivalent at APPS training scale. Practitioners should invest in execution infrastructure, not reward function engineering. This is a replicable, controlled negative result.

\textbf{The APPS coverage paradox is a new empirical observation.} 85.32\% competition-level solution coverage despite 0.0 SFT pass@1 is a striking decoupling of dataset quality from evaluation performance. Future code LLM training work should not assume dataset coverage improvements automatically translate to hard-benchmark gains.

\subsection{Limitations}

\textbf{L1: Smoke-scale evaluation.} All $\Delta$ values are proxies from 500 APPS samples, 62 GRPO steps, N=50 evaluation per benchmark. RLEF checkpoint was not saved (process timeout). Full-scale evaluation is deferred to Phase 5. \textit{Why acceptable:} MUST\_WORK gate confirmed mechanism; JT confirms directional claim; infrastructure is validated.

\textbf{L2: Mechanism unverified.} The partial-success gradient mechanism was not directly observed. h-m2 reward monitoring used max\_new\_tokens=128 (truncation artifact producing 0.0 reward). Re-test with max\_new\_tokens$\geq$512 is required. \textit{Why acceptable:} Empirical outcome (difficulty-scaling advantage) confirmed independently of mechanism.

\textbf{L3: Reward comparison uses proxy estimates.} Full RLEF-Binary training was not completed. Null result direction is consistent with two independent papers but requires full verification. \textit{Why acceptable:} Directional evidence is literature-consistent; full comparison is planned.

\textbf{L4: Statistical values conditional on proxy data.} JT z=+56.10 and bootstrap p-values reflect bootstrap construction from point estimates, not independent experimental replications. \textit{Why acceptable:} Methods are valid conditional on inputs; full replication is Phase 5 objective.

\textbf{L5: Evaluation scope limited.} Function-level Python correctness-only evaluation. Repository-level tasks (SWE-bench), multi-language, and process reward metrics are not tested.

\subsection{Broader Impact}

\textbf{Positive:} Open-source pipeline enables community replication and extension. Negative result on reward formulation saves practitioners from unnecessary engineering complexity. Reproducibility contribution closes a gap in the RLEF literature.

\textbf{Potential concerns:} Improved code generation capability could accelerate automated code production without sufficient quality control. Our evaluation measures correctness only --- not security, maintainability, or broader software engineering properties. Sandbox isolation is critical for safe deployment of any code execution pipeline; our implementation uses subprocess isolation with resource limits.
